% Supplementary result: a bounded-coefficient obstruction to sparse sampling.
% Prepared independently from the repository's one-sided prime-sum mechanism.
% Requires amsmath, amssymb, amsthm; theorem/proof environments provided by parent.

\section{A bounded-coefficient obstruction to sparse sampling}
\label{sec:sampling-obstruction}

The one-sided analytic criterion in \texttt{grh-lean} applies to a
positively weighted, smoothed arithmetic sum.  Positivity also permits
interpolation between sampled values.  The fourth-power sampling principle
itself is already in Setiawan's November 2026 preprint~\cite{Setiawan2026}: its Sections~3 and~4
use convexity to recover a continuous bound from
suitably spaced samples.  The same preprint includes single-interval
sharpness for convex interpolation.  We do not claim those statements as
new.  The theorem below supplies a different strengthening of the
obstruction: one globally defined sequence, bounded positive coefficients,
exactly vanishing samples at \emph{every} fifth power, and a genuine Mellin
pole that all those samples miss.

\begin{theorem}[Bounded positive coefficients with an invisible pole]
\label{thm:invisible-pole}
There is an explicitly defined sequence of real numbers $(a_n)_{n\geq1}$
with
\[
                         0<a_n<4
\]
such that the function
\[
 R(x)=\sum_{n\leq x}a_n\left(\frac1n-\frac1x\right)-\log x,
                    \qquad x\geq1,
\]
has all the following properties.
\begin{enumerate}
\item $R(m^5)=0$ for every positive integer $m$.
\item For every $x\geq1$,
\[
0\leq R(x)-\frac12 x^{-2/5}\sin^2(\pi x^{1/5})
                       \leq \frac1{2x^2}.
\tag{S1}\label{eq:sampling-error}
\]
In particular $R\geq0$, $R(x)=O(x^{-2/5})$, and, for every
$\beta<3/5$, the upper estimate $R(x)=O_+(x^{\beta-1})$ is false.
\item The Mellin transform
\[
                  \Phi(s)=\int_1^\infty R(x)x^{-s}\,dx
\]
is absolutely convergent for $\Re s>3/5$.  It has a meromorphic
continuation to $\Re s>-1$ whose only pole is a simple pole at $s=3/5$,
with residue $1/4$.
\end{enumerate}
The notation $O_+$ means an upper bound with a fixed positive constant;
no lower bound is intended.  The coefficients are real numbers indexed
by integers, not integer-valued coefficients.
\end{theorem}

\begin{proof}
Define, for $x\geq1$,
\[
 h(x)=x^{3/5}\sin^2(\pi x^{1/5}),\qquad
 F(x)=x\log x+\frac12h(x).
\]
Thus $F(1)=0$.  Set
\[
 a_1=F(2),\qquad
 a_n=n\bigl(F(n+1)-2F(n)+F(n-1)\bigr)\quad(n\geq2).
\tag{S2}\label{eq:explicit-coefficients}
\]
These formulas define every coefficient without a limiting choice.

We first verify positivity and the uniform bound.  Differentiation gives
\begin{align*}
 xh''(x)={}&-\frac6{25}x^{-2/5}\sin^2(\pi x^{1/5})\\
 &+\frac{2\pi}{25}x^{-1/5}\sin(2\pi x^{1/5})
   +\frac{2\pi^2}{25}\cos(2\pi x^{1/5}).
\end{align*}
Consequently, using $x\geq1$ and $\pi<4$,
\[
 |xh''(x)|\leq\frac{6+2\pi+2\pi^2}{25}<\frac{46}{25}.
\]
It follows that
\[
              \frac{2}{25x}<F''(x)<\frac{48}{25x}<\frac2x.
\tag{S3}\label{eq:convexity-bounds}
\]
For $n\geq2$, the standard second-difference identity
\[
 F(n+1)-2F(n)+F(n-1)
   =\int_{-1}^{1}(1-|u|)F''(n+u)\,du
\]
now proves
\[
            0<a_n<\frac{2n}{n-1}\leq4.
\]
Also $0<a_1=2\log2+\tfrac12 2^{3/5}\sin^2(\pi2^{1/5})<3$.

Let $IF$ be the continuous function that agrees with $F$ at every
positive integer and is affine between consecutive integers.  Its slope
on $[m,m+1]$ is $F(m+1)-F(m)$.  The function
\[
                  Q(x)=\sum_{n\leq x}a_n(x/n-1)
\]
has the same initial value $Q(1)=0$ and, on that interval, its slope is
\begin{align*}
 \sum_{n\leq m}\frac{a_n}{n}
 &=F(2)+\sum_{n=2}^{m}\bigl(F(n+1)-2F(n)+F(n-1)\bigr)\\
 &=F(m+1)-F(m).
\end{align*}
Thus $Q=IF$ on $[1,\infty)$ and
\[
 R(x)=\frac{IF(x)}x-\log x
     =\frac12x^{-2/5}\sin^2(\pi x^{1/5})+\varepsilon(x),
 \qquad
 \varepsilon(x)=\frac{IF(x)-F(x)}x.
\tag{S4}\label{eq:decomposition}
\]
Since $F$ is convex, $\varepsilon\geq0$.  On $[m,m+1]$,
\eqref{eq:convexity-bounds} and the elementary chord-error bound give
\[
0\leq IF(x)-F(x)
 \leq\frac{(x-m)(m+1-x)}m\leq\frac1{4m}.
\]
Dividing by $x$, and using $x\leq m+1\leq2m$, gives
$0\leq\varepsilon(x)\leq1/(2x^2)$.  At every integer $x$,
$\varepsilon(x)=0$.  In particular, $x=m^5$ makes both terms in
\eqref{eq:decomposition} zero, proving the exact sample identities.

At $x_m=(m+\tfrac12)^5$, however, the sine-squared term equals one, so
\[
             \frac{R(x_m)}{x_m^{\beta-1}}
                    \geq\frac12 x_m^{3/5-\beta}\longrightarrow\infty
                    \qquad(\beta<3/5).
\]
This proves the failure of the proposed continuous upper bound.

It remains to prove the claimed analytic continuation, rather than merely
infer a singularity from the size of the intervening humps.  By
\eqref{eq:sampling-error},
\[
                 H(s)=\int_1^\infty\varepsilon(x)x^{-s}\,dx
\]
is holomorphic in $\Re s>-1$, by absolute convergence uniformly on compact
subsets (the same estimate permits differentiation under the integral).
The identity $\sin^2t=(1-\cos2t)/2$ and the substitution $x=y^5$ give,
initially for $\Re s>3/5$,
\[
 \Phi(s)=\frac1{4(s-3/5)}
          -\frac54 I(2-5s)+H(s),
 \qquad
 I(z)=\int_1^\infty y^z\cos(2\pi y)\,dy.
\tag{S5}\label{eq:pole-decomposition}
\]
Initially $I$ is absolutely convergent for $\Re z<-1$.  Two integrations
by parts give there
\[
               I(z)=-\frac{z}{4\pi^2}
                     -\frac{z(z-1)}{4\pi^2}I(z-2).
\tag{S6}\label{eq:oscillatory-continuation}
\]
The boundary terms at infinity vanish in this half-plane, and at $1$ we
used $\sin(2\pi)=0$ and $\cos(2\pi)=1$.  Repeatedly applying
\eqref{eq:oscillatory-continuation} continues $I$ holomorphically to every
half-plane $\Re z<2N-1$.  There are no variable denominators, so these
continuations are entire and agree on overlaps by the identity theorem.
Hence the last two terms in \eqref{eq:pole-decomposition} are holomorphic
on $\Re s>-1$.  The first term is the asserted unique simple pole and has
residue $1/4$.
\end{proof}

\paragraph{What the obstruction establishes.}
For any $1/2<\beta<3/5$, all the sample bounds
$R(m^5)\leq m^{5(\beta-1)}$ hold, even with a zero left-hand side, while
the Mellin transform has a positive real residue at a pole to the right
of $\beta$.  Thus positivity, uniformly bounded coefficients, exact
fifth-power observations, and meromorphic continuation across the
relevant half-plane do not suffice to reconstruct the one-sided estimate
needed by a Landau argument.  The actual von Mangoldt weights have more
arithmetic structure; this example does not rule out exploiting that
structure in a sparser criterion.  The continuation proved here is to
$\Re s>-1$.  We do not assert that this construction satisfies the exact
whole-plane \texttt{IsPolar} definition in the Lean repository.

\paragraph{A family of examples.}
The same construction works with every integer $p\geq4$, replacing
$3/5$ by $1-2/p$, $x^{1/5}$ by $x^{1/p}$, and fifth powers by $p$th
powers.  The pole lies at $1-2/p$ and still has residue $1/4$.
Indeed, writing $r=1/p$, the analogue of the second-derivative bound is
\[
 |xh_p''(x)|\leq 2r(1-2r)+\pi r(1-3r)+2\pi^2r^2
 =C_p:=\frac{2+\pi}{p}+\frac{2\pi^2-3\pi-4}{p^2}.
\]
Both coefficients in the last expression are positive, so $C_p$ decreases
for $p\geq4$.  Moreover
$C_4=1/4+\pi/16+\pi^2/8<2$, for example using $\pi<22/7$.
Thus $0<F_p''(x)<2/x$ and the same bounds on the weights follow.
The interpolation error remains $O(x^{-2})$, and the oscillatory Mellin
term is $-(p/4)I(p(1-s)-3)$, which is entire in $s$.
